\section{Related Work}
\label{sec:related_work}
\subsection{Label-Free Reinforcement Learning}

A significant trend in recent research is Label-Free Reinforcement Learning, which aims to improve LLM reasoning without human-annotated data. Many such methods use the model's own outputs as a reward signal. This includes leveraging sequence-level confidence~\citep{Li2025ConfidenceIA, Prabhudesai2025MaximizingCA}, the consistency of answers derived from varied reasoning paths~\citep{Zhang2025ConsistentPL, Zuo2025TTRLTR, Zhang2025RightQI,zhou2025evolving,prasad2024self}, minimizing the output entropy~\citep{Agarwal2025TheUE, Cheng2025ReasoningWE}, or even random~\citep{shao2025spuriousrewardsrethinkingtraining} or negative reward~\citep{Zhu2025TheSE}. These signals are often used within self-training loops where models fine-tune on their own most plausible solutions~\citep{Shafayat2025CanLR, Zhao2025LearningTR}. While these methods all rely on a pre-existing set of unlabeled problems, \method{} removes the need for any seed dataset.
\subsection{Self-Play in Large Language Models}
\label{sec:related_work_self_play}

The paradigm of self-play, where models take on dual roles to create a self-improvement loop~\citep{chen2024self,zhang2024survey}, has recently been adapted to improve language models without human data. This approach has been particularly fruitful in verifiable domains like code generation, where a ``Coder'' agent's program is verified by a ``Tester'' agent's unit tests~\citep{Lin2025LearningTS, Wang2025CoEvolvingLC, Pourcel2025SelfImprovingLM,jiang2025verse}. More advanced frameworks push autonomy further by learning to generate the problems themselves, creating an adaptive curriculum from a small seed of examples or from scratch~\citep{Zhao2025AbsoluteZR, li2025videohalluevaluatingmitigatingmultimodal,Zhou2025SelfChallengingLM, Fang2025SeRLSR}. Our work distinguishes itself by extending this paradigm to general reasoning domains that lack such verifiable environments, like coding tasks.

\subsection{Reinforcement Learning with Verifiable Rewards (RLVR)} Reinforcement Learning with Verifiable Rewards has been widely adopted as a versatile paradigm for enhancing LLMs across a multitude of tasks~\citep{li2025self,DeepSeekAI2025DeepSeekR1IR,Shao2024DeepSeekMathPT}. Its effectiveness is demonstrated in diverse applications such as relation extraction \citep{dai2025r1}, interactive GUI navigation \citep{Shi2025MobileGUIRLAM} and search-engine utilization \citep{Jin2025SearchR1TL}. While early implementations relied on rule-based verifiers, recent work has begun to explore more sophisticated, model-based verifiers~\citep{Ma2025GeneralReasonerAL,li2025semanticallyawarerewardsopenendedr1,li2024fg}.